\section{teorema di Stone Weierstrass}

\begin{theorem}[Stone Weierstrass]
Sia $X$ uno spazio metrico compatto.
Sia ${\mathcal A} \subset C(X)$ un'algebra 
cioè un sottospazio vettoriale dello spazio delle funzioni continue 
su $X$ a valori in $\RR$
tale che se $f,g\in {\mathcal A}$ allora anche $f\cdot g\in {\mathcal A}$.
Supponiamo che
\begin{enumerate}
\item $1 \in {\mathcal A}$ ($1$ è la funzione costante $1(x)=1$);
\item ${\mathcal A}$ separa i punti di $X$, cioè
\[
\forall x \neq y \in X \quad \exists f \in {\mathcal A} \text{ tale che } f(x)\neq f(y).
\]
\end{enumerate}

Allora ${\mathcal A}$ è densa in $C(X)$ rispetto alla norma uniforme,
cioè 
\[
\forall \eps>0\, \forall f\in C(X)\,\exists g\in {\mathcal A}\colon \Abs{f-g}_\infty < \eps.
\]
\end{theorem}

Per la dimostrazione ci servono alcuni lemmi intermedi.

\begin{lemma}
Siano $X$ e ${\mathcal A}$ come nelle ipotesi del teorema.
Dato $x_0\in X$, dato $U$ intorno aperto di $x_0$
esiste $V$ intorno di $x_0$, $V\subset U$
tale che per ogni $\eps>0$ 
esiste $f\in {\mathcal A}$ tale che 
\begin{enumerate}
\item    $0 \le f \le 1$,
\item    $f(x) < \eps$ per ogni $x\in V$,
\item    $f(x) > 1-\eps$ per ogni $x \in X\setminus U$.
\end{enumerate}
\end{lemma}
%
\begin{proof}
Per ogni $x\in X\setminus U$ 
consideriamo $g_x \in {\mathcal A}$ 
tale che $g_x(x)\neq g_x(x_0)$.
Definiamo $h_x = g_x - g_x(x_0)\cdot 1 \in {\mathcal A}$.
Così $h_x(x) \neq  h_x(x_0) = 0$.
Definiamo
\[
 p_x = \frac{h_x^2}{\Abs{h_x}_\infty} \in {\mathcal A}
\]
così $p_x(x_0)=0$, $p_x(x)>0$,
$0\le p_x \le 1$.

L'insieme 
\[
U_x = \ENCLOSE{y\in X\colon p_x(y)>0}
\]
è un intorno aperto di $x$
dunque $X\setminus U$ è un compatto 
ricoperto dagli $U_x$ al variare di $x\in X\setminus U$.
Posso estrarre un sottoricoprimento finito:
dunque esistono $x_1,\dots,x_n\in X\setminus U$ 
tali che $X\setminus U \subset U_{x_1} \cup \dots \cup U_{x_m}$.
Poniamo
\[
  p = \frac{p_{x_1}+ \dots + p_{x_m}}{m} \in {\mathcal A}
\]
così da avere $0\le p \le 1$, $p(x_0)=0$,
$p(x)>0$ per ogni $x\in X\setminus U$.
Visto che $X\setminus U$ è compatto 
esiste $\delta = \min p(X\setminus U)>0$ 
e dunque $p(x)\ge \delta > 0$ e $0<\delta\le 1$.
Sia 
\[
  V = \ENCLOSE{x\in X \colon p(x) < \frac \delta 2}.
\]
Visto che $p(x_0)=0$ si ha che $V$ è un intorno di $x_0$,
e visto che 
$p(x) \ge \delta$ su $X\setminus U$
risulta $V\subset U$.
Sia $k\in \NN$ tale che 
\[
 \frac 1 \delta < k \le \frac 1 \delta + 1 
 = \frac{\delta+1}{\delta}
 < \frac 2 \delta
\]
così $1 < k \delta < 2$.
Poniamo
\[
 q_n(x) = \Enclose{1-p^n(x)}^{k^n} \in {\mathcal A}
\]
così $0\le q_n \le 1$, $q_n(x_0) = 1$.
Ricordando la disuguaglianza di Bernoulli
\[
  (1+t)^n \ge 1 + nt, \qquad \forall t\ge -1
\]
per ogni $x\in V$ si ha $k p(x)\le k \frac \delta 2 < 1$ 
dunque
\[
q_n(x) \ge 1- k^n p^n(x)
\ge 1 - \enclose{k \frac \delta 2}^n \to 1
\]
ovvero $q_n \to 1$ uniformemente su $V$.
Mentre se $x\in X\setminus U$ 
si ha $kp \ge k\delta > 1$
\begin{align*}
 q_n(x) 
 &= \Enclose{1-p^n(x)}^{k^n}
 \le \frac{1+k^np^n(x)}{k^np^n(x)}
 \Enclose{1-p^n(x)}^{k^n}\\
 &\le \frac{\Enclose{1+p^n(x)}^{k^n}}{k^n p^n(x)}
 \Enclose{1-p^n(x)}^{k^n}\\
 &= \frac{\Enclose{1-p^{2n}(x)}^{k^n}}{k^n p^n(x)}
 \le \frac{1}{k^n p^n(x)}
 \le \frac{1}{(k\delta)^n} \to 0.
\end{align*}
Dunque $q_n\to 0$ uniformemente su $X\setminus U$.

Prendendo $n$ sufficientemente grande 
si ha $q_n<\eps$ su $X\setminus U$ 
e $q_n > 1-\eps$ su $V$.

La funzione $f=1-q_n$ soddisfa le proprietà richieste.
\end{proof}

\begin{lemma}
Siano $X$ e ${\mathcal A}$ come nel teorema.
Siano $A,B\subset X$ compatti disgiunti.
Per ogni $\eps>0$ esiste $f\in {\mathcal A}$ tale che 
\begin{enumerate}
\item    $0 \le f \le 1$,
\item    $f(x) < \eps$ per ogni  $x\in A$,
\item    $f(x) > 1-\eps$ per ogni $x \in B$.
\end{enumerate}
\end{lemma}
%
\begin{proof}
Sia $U\supset A$ aperto disgiunto da $B$.
Per ogni $x\in A$ esiste $V_x\subset U$ 
un intorno aperto di $x$ come nel lemma precedente.
Gli aperti $V_x$ ricoprono il compatto $A$ 
dunque esistono $x_1, \dots, x_m\in A$ tali che 
$V_{x_1} \cup \dots \cup V_{x_m} \supset A$
ed esistono $g_1, \dots, g_m \in {\mathcal A}$ 
tali che $0\le g_k \le 1$, $g_i < \frac \eps m$ su $V_{x_i}$
e $g_i > 1 - \frac \eps m$ su $X\setminus U$.

Posto $g= g_1 \cdot g_2 \cdots g_m$ 
si ha $0\le g \le 1$ e per ogni $x\in A$ 
esiste $x_i$ tale che $x\in V_{x_i}$.
Visto che $g_i(x) < \frac \eps m$ si ha 
$g(x) < \frac \eps m$.
Inoltre per ogni $x\in B$, $x \not \in \bigcup V_{x_i}$
e per ogni $i$ si ha $g_i(x) < 1 - \frac \eps m$.
Dunque 
\[
g > \enclose{1-\frac\eps m}^m \ge 1 - m \frac \eps m = 1-\eps.
\]
\end{proof}

\begin{proof}[Dimostrazione del teorema.]
Sia $f\in C(X)$ e $\eps>0$.
Voglio mostrare che esiste $g\in {\mathcal A}$ tale che
$\abs{f-g}_\infty \le \eps$.
Eventualmente considerando $f + \Abs{f}_\infty$ posso 
supporre $f\ge 0$.
Eventualmente considerando $\frac {f}{3\eps}$
posso supporre $\eps = 3$.

Prendiamo $N\in \NN$, $N \ge \Abs{f}_\infty+1$.
Per $j=1, \dots, N$ definiamo 
\begin{align*}
    A_j &= \ENCLOSE{x\in X\colon f(x) \le j}\\
    B_j &= \ENCLOSE{x\in X\colon f(x) \ge j+1}
\end{align*}
così $A_j$ e $B_j$ sono compatti disgiunti.

Per il lemma precedente esiste
$g_j\in{\mathcal A}$, $0\le g_j \le 1$ tale che
$0 \le g_j \le \tfrac 1 N$
su $A_j$ 
e $1 - \tfrac 1 N \le g_j \le 1$ su $B_j$.

Si definisce
\[
g = \sum_{j=1}^N g_j(x).
\]
Per ogni $x\in X$ esiste $m\in \NN$, $0\le m\le N$, 
tale che $m\le f(x) \le m+1$. 
Dunque 
$x \in B_j$ per $j=1, \dots, m-1$
e $x \in A_j$ per $j=m+1, \dots N$.

Dunque si ha  
\[
  (m-1)\cdot (1-\tfrac 1 N) + (N-m+1)\cdot 0 
  \le g(x) 
  \le m\cdot 1 + (N-m)\cdot \frac 1 N
\]
e in particolare $m-2 \le g(x) \le m+1$.
In conclusione 
\[
m - (m+1)
\le f(x) - g(x) 
\le m+1 - (m-2)
\]
ovvero $-1 \le f(x) - g(x) \le 3$,
da cui $\abs{f-g}_\infty \le 3$.
\end{proof}